\documentclass[25pt, a0paper, portrait]{tikzposter}
\usetheme{Rays}
\usecolorstyle{Default}
\usepackage[utf8]{inputenc}
\usepackage[T1]{fontenc}
\usepackage[french]{babel}
\usepackage{amsmath,amssymb}

\title{Makeitalive : Pipeline de Données et Analyse des Résultats}
\author{Arthur et Camille}
\institute{Projet IA Générative - Dataset et Performances}

\begin{document}
\maketitle

\begin{columns}
    \column{0.5}
    \block{1. Stratégie d'Extraction}{
        Le pipeline automatise la création de paires d'images $(I\_A, I\_B)$ :
        \begin{itemize}
            \item \textbf{Source :} Vidéos YouTube 4K (Landscapes).
            \item \textbf{Sampling :} Extraction toutes les 5 secondes.
            \item \textbf{Frame Gap :} Écart de 3 frames pour un mouvement fluide.
        \end{itemize}
    }

    \block{2. Filtrage Intelligent des Données}{
        Éviter les "cuts" via deux métriques robustes :
        \begin{itemize}
            \item \textbf{Hist Distance} : Bhattacharyya ($>0.35$).
            \item \textbf{Thumbnail MAE} : Différence structurelle ($>0.25$).
        \end{itemize}
        Ce filtrage assure que le modèle n'apprend que de vrais mouvements et non des sauts de montage.
    }

    \block{3. Paramètres de l'Architecture}{
        \begin{itemize}
            \item \textbf{Input :} $512 \times 512$ (Center Crop).
            \item \textbf{Optimiseur :} Adam ($10^{-4}$).
            \item \textbf{Loss :} Mean Squared Error (MSE) sur le flow.
            \item \textbf{Epochs :} 100 sur dataset local dédié.
        \end{itemize}
    }

    \column{0.5}
    \block{4. Analyse Qualitative}{
        \begin{itemize}
            \item \textbf{Points forts :} Excellence sur l'eau et les reflets.
            \item \textbf{Limites :} Quelques artefacts sur les nuages denses à cause de l'ambiguïté du flow.
            \item \textbf{Bords :} \texttt{mode='reflect'} indispensable pour l'esthétique.
        \end{itemize}
    }

    \block{5. Performances (Convergence de la Loss)}{
        \begin{center}
        \begin{tabular}{|c|c|c|}
            \hline
            \textbf{Epoch} & \textbf{Training Loss} & \textbf{Validation Loss} \\
            \hline
            1 & 0.452 & 0.480 \\
            10 & 0.120 & 0.145 \\
            50 & 0.045 & 0.052 \\
            100 & 0.038 & 0.041 \\
            \hline
        \end{tabular}
        \end{center}
    }

    \block{6. Perspectives}{
        \begin{itemize}
            \item \textbf{Segmentation sémantique} : Appliquer des flows spécifiques au ciel vs l'eau.
            \item \textbf{Inférence multi-échelle} : Améliorer la précision des détails.
            \item \textbf{Raffinement Temporel} : Prédiction auto-répressive pour des loops fluides.
        \end{itemize}
    }

    \block{7. Impact}{
        Le projet démontre qu'une architecture légère peut produire des animations qualitatives avec très peu de données d'entraînement.
    }
\end{columns}

\end{document}
